For $a,b,c \in \NN$ we denote by $\langle a,b,c\rangle \in \FF^{ab} \otimes \FF^{bc} \otimes \FF^{ca}$ the %
matrix multiplication tensor $\langle a,b,c \rangle = \sum_{i=1}^a \sum_{j=1}^b \sum_{k=1}^c e_{i,j} \otimes e_{j,k} \otimes e_{k,i}$. 
These tensors have the multiplicative property $\langle a,b,c\rangle \otimes \langle x,y,z\rangle = \langle ax, by, cz\rangle$. They are naturally related to the parameters~$\subrank_i$ as follows:

\begin{lemma}\label{lem:Qi-mamu}
Let $T$ be a tensor and $r \in \NN$. The following holds:
\begin{itemize}
\item $\subrank_1(T) \geq r$ if and only if $T \geq \langle 1, 1,r \rangle$.
\item $\subrank_2(T) \geq r$ if and only if $T \geq \langle r, 1,1 \rangle$.
\item $\subrank_3(T) \geq r$ if and only if $T \geq \langle 1, r, 1 \rangle$.  
\end{itemize}
\end{lemma}

\begin{lemma}\label{lem:Qi-square}
    If there are distinct ${\ell_1},{\ell_2}$ such that $\subrank_{\ell_1}(T), \subrank_{\ell_2}(T) \geq r$, then $\subrank(T^{\boxtimes 2}) \geq r$.
\end{lemma}

\begin{proof}
Suppose $\subrank_1(T), \subrank_2(T) \geq r$ (the proof for the other cases goes the same). Then we have that
$
T \geq \langle1, 1, r\rangle =  \sum_{i=1}^r e_1 \otimes e_i \otimes e_i$ and 
$T \geq \langle r, 1, 1\rangle = \sum_{i=1}^r e_i \otimes e_1 \otimes e_i$.
%
We multiply these inequalities to get
$T^{\boxtimes 2} \geq \langle r,1,r\rangle= \sum_{i,j=1}^r e_i \otimes e_j \otimes (e_i \otimes e_j).
$
We apply the projection $(e_i \otimes e_j) \mapsto \delta_{i=j} e_i$ to the third tensor leg to get
$
T^{\boxtimes 2} \geq \sum_{i=1}^r e_i \otimes e_i \otimes e_i
$
which proves the claim.
\end{proof}

We are now ready to prove the fourth-root lower bound on the asymptotic subrank. After that we will prove the stronger cube-root lower bound (\autoref{th:cube-root-lb}).

\begin{theorem}\label{th:fourth-root}
    Let $T \in \FF^{n_1} \otimes \FF^{n_2} \otimes \FF^{n_3}$ be concise. Then $\asympsubrank(T) \geq (\min_\ell n_\ell)^{1/4}$.
\end{theorem}
\begin{proof}
The proof is an application of \autoref{th:uncertainty-principle-n}.
That theorem only holds over large enough fields however. 
Since the asymptotic subrank does not change under field extension \cite[Theorem~3.10]{strassen1988asymptotic}, we may assume that our base field $\FF$ is infinite (by working over the algebraic closure of the original field) so that the field size condition of \autoref{th:uncertainty-principle-n} is satisfied. %
We use that there are distinct $\ell_1, \ell_2 \in [3]$ such that 
%\refnote{13}{missing square root on $n_i$}\ILnote{maybe it was missing in an earlier version?}\Jshort{yes we fixed this earlier} 
$\subrank_{\ell_1}(T), \subrank_{\ell_2}(T) \geq \sqrt{\min_\ell n_\ell}$ (\autoref{cor:twoLargeRankDirs}).
Then $\asympsubrank(T^{\boxtimes 2}) \geq \sqrt{\min_\ell n_\ell}$ (\autoref{lem:Qi-square}) which gives the claim. %
\end{proof}




%

It turns out that the above construction is not optimal. We will now give a slightly more elaborate construction that leads to the better cube-root bound:

\begin{theorem}\label{th:cube-root-lb}
    Let $T \in \FF^{n_1} \otimes \FF^{n_2} \otimes \FF^{n_3}$ be concise. Then $\asympsubrank(T) \geq (\min_{\ell} n_\ell)^{1/3}$.
\end{theorem}

The proof of \autoref{th:cube-root-lb} makes use of basic properties of matrix multiplication tensors which we discuss in the following two lemmas. 
%\refnote{14}{The following is a repetition of matrix multiplication notation description -- it is used above in this page already} %Recall that we use the notation~$\langle a,b,c\rangle$ for~$a,b,c \in \NN$ to denote the matrix multiplication tensors.
%\Jshort{Removed it.}

\begin{lemma}\label{lem:cube-mamu}
    $T^{\boxtimes 3} \geq \langle \subrank_2(T),\subrank_3(T),\subrank_1(T) \rangle.$
\end{lemma}

\begin{proof}
    From \autoref{lem:Qi-mamu} we have %three inequalities
    %
    %
    %
    %
    %
    %
    $T \geq \langle \subrank_2(T),1,1\rangle$, 
    $T \geq \langle 1,\subrank_3(T),1\rangle$, and  
    $T \geq \langle 1,1,\subrank_1(T)\rangle$.
    %
We multiply these inequalities using the Kronecker product to get the claim.
\end{proof}


\begin{lemma}\label{lem:asympsub-min}
    $\asympsubrank(T^{\boxtimes 3}) \geq \min_{{\ell_1},{\ell_2}} \subrank_{\ell_1}(T) \subrank_{\ell_2}(T).$
\end{lemma}
\begin{proof}
    It is known that for every $a,b,c \in \NN$ the asymptotic subrank of the matrix multiplication tensor $\langle a,b,c\rangle$ is given by $\asympsubrank(\langle a,b,c\rangle) = \min \{ab, bc, ac\}$ \cite[Equation~4.6]{strassen1988asymptotic}. We combine this with \autoref{lem:cube-mamu} to get the claim.
\end{proof}

\begin{proof}[Proof of \autoref{th:cube-root-lb}]
As in the proof of \autoref{th:fourth-root} we may assume that the base field $\FF$ is infinite.
%
%
%
We use $\asympsubrank(T^{\boxtimes 3}) \geq \min_{{\ell_1},{\ell_2}} \subrank_{\ell_1}(T) \subrank_{\ell_2}(T)$ (\autoref{lem:asympsub-min}) and we use that for every distinct ${\ell_1}, {\ell_2}, {\ell_3} \in [3]$ we have $\subrank_{\ell_1}(T)\subrank_{\ell_2}(T) \geq n_{\ell_3}$ (\autoref{th:uncertainty-principle-n}) to obtain
$\asympsubrank(T^{\boxtimes 3}) \geq \min_{{\ell_1},{\ell_2}} \subrank_{\ell_1}(T) \subrank_{\ell_2}(T) \geq \min_{\ell_3} n_{\ell_3}$. %
\end{proof}
